% Vector figures for the evidence-led IndiCare report.
% Shared palette/styles and system views are provided by diagram_system.tex.
% Source-audited deployment, component, class, data and security maps.
% Shared palette/geometry lives in diagram_styles.tex. Units: cm.
\newcommand{\ArchitectureDiagram}{%
\begin{figure}[H]\centering
\begin{tikzpicture}
% Three physical responsibilities; colour separates runtime from offline work.
\path[fill=bluewash,rounded corners=3pt] (11.5,.85) rectangle (15.6,-8.35);
\node[dHeading,anchor=west] at (11.75,.55) {SELF-HOSTED PC};
\node[dBox,text width=2.8cm,minimum height=3.25cm] (wearable) at (1.6,-1.5)
 {\textbf{01 / Wearable}\\[4pt]KEYES + ESP32\\MAX30102 / GY906\\MPU6050\\[4pt]Acquire + flag validity\\Buzzer output: GPIO25};
% Phone silhouette remains legible without using a mock screenshot.
\draw[draw=accent,fill=white,line width=.85pt,rounded corners=6pt] (5.25,.25) rectangle (8.75,-3.25);
\draw[stroke,line width=1.3pt,line cap=round] (6.55,.07)--(7.45,.07);
\node[dNote,text=ink,text width=2.95cm,anchor=north west] at (5.52,-.2)
 {\textbf{02 / Adult Android app}\\[5pt]CareMonitor\\Authenticate wearable\\Encrypt upload queue\\Readings + history\\Saved Learning insights};
\coordinate (phoneIn) at (5.25,-1.35);\coordinate (phoneOut) at (8.75,-1.35);
\node[dBox,text width=2.9cm,minimum height=2.65cm] (api) at (13.55,-1.65)
 {\textbf{03 / FastAPI}\\[4pt]REST JSON /api/v1\\Roles + approved links\\Readings / history / inbox\\SQLite + encrypted fields};
\draw[dFlow] (wearable.east|-phoneIn)--(phoneIn);
\draw[dReturn] (5.25,-1.85)--(3.22,-1.85);
\node[dLabel,text width=1.8cm] at (4.23,-.6) {Classic SPP\\NDJSON telemetry};
\node[dLabel,text width=1.8cm] at (4.23,-2.55) {CM1 mutual\\authentication};
\draw[dFlow] (phoneOut)--(api.west|-phoneOut);
\draw[dReturn] (api.west|-0,-1.85)--(8.75,-1.85);
\node[dLabel,text width=2.1cm] at (10.15,-.65) {HTTPS / TLS\\REST JSON + token};
\node[dLabel,text width=2.1cm] at (10.15,-2.55) {Upload / query\\Parsed response};
% Caregiver has a real connection, rather than an unconnected paragraph.
\node[dData,text width=2.8cm,minimum height=1.6cm] (care) at (1.6,-4.7)
 {\textbf{04 / Caregiver app}\\Approved adult selection\\History + private notices};
\draw[draw=accent,line width=.85pt,{Stealth[length=2mm]}-{Stealth[length=2mm]}]
 (care.east)--(10.9,-4.7)--(10.9,-3.05)--(api.south west);
\node[dLabel,text width=5.7cm] at (7.1,-4.2) {HTTPS long polling and authorised GET\\Adult approval controls the shared resource};
% Independent training is separate from deployed native inference.
\node[dBox,fill=light,text width=2.8cm,minimum height=1.7cm] (csv) at (1.6,-7.25)
 {\textbf{Synthetic dataset}\\Grouped row split\\Educational labels};
\node[dBox,fill=light,text width=3cm,minimum height=1.7cm] (train) at (7,-7.25)
 {\textbf{Offline model comparison}\\Logistic Regression / RF\\Select on validation};
\node[dData,text width=2.9cm,minimum height=1.7cm] (model) at (13.55,-7.25)
 {\textbf{Loaded health model}\\Hash + version checks\\Class + confidence score};
\draw[dReturn] (csv)--node[dLabel,above]{Train / evaluate}(train);
\draw[dReturn] (train)--node[dLabel,above]{Saved artefact}(model);
\draw[dFlow] (api.south)--node[dLabel,right,text width=1.9cm]{Synthetic-only\\Learning API}(model.north);
\node[dNote,anchor=west] at (.05,-9.05) {\textbf{Legend}\quad Solid: implemented runtime communication.\quad Dashed: reply or offline artefact path, as labelled.};
\end{tikzpicture}
\caption{Complete system and its separate care and synthetic-learning paths. Physical SPP, API-backed native views and saved synthetic confidence were exercised in the controlled code-22 screen campaign; model training was evaluated independently. USB forwarding carried HTTPS in that campaign; reachable LAN deployment uses the same API. The two roles were exercised successively on one S22.}\label{fig:architecture}
\end{figure}}

\newcommand{\ProtocolStackFigure}{%
\begin{table}[H]\centering\small
\caption{Implemented communication boundaries. Cryptographic controls and local persistence are distinguished from transport protocols.}\label{tab:protocol-stack}
\begin{tabularx}{\textwidth}{@{}P{28mm}P{47mm}Y@{}}\toprule
Boundary & Protocol / format & Contract and protection\\\midrule
Sensors / ESP32 & I\textsuperscript{2}C / SMBus, 100 kHz & GPIO21 SDA, GPIO22 SCL; 7-bit addresses 0x57 / 0x5A / 0x68. Common ground and compatible logic.\\
ESP32 / phone & Classic Bluetooth SPP; RFCOMM UUID 0x1101; UTF-8 NDJSON & Encrypted link plus CM1 HMAC-SHA256 proofs. One active phone; telemetry after CM1 READY. This firmware does not expose BLE GATT.\\
Phone / PC carrier & USB adb reverse, or LAN IP/TCP & USB maps phone loopback to PC loopback. LAN normally uses port 8443. Both carry the same certificate-verified HTTPS application traffic.\\
App / FastAPI & HTTPS REST JSON, /api/v1 & Certificate chain and hostname checks; Bearer token; per-resource role and approved-link authorisation.\\
Private care inbox & Authenticated HTTPS long polling & Incremental durable inbox IDs. Delivery acknowledgement, human review and episode recovery are separate operations.\\
Services / SQLite & sqlite3 transactions and WAL & Unique sample IDs, atomic fanout and encrypted sensitive fields. AES-GCM is field protection, not full-file database encryption.\\
Optional board clock & Wi-Fi / IP / UDP / NTP & Time synchronisation requires reachability. Missing wearable timestamps use phone receipt time; NTP does not authenticate measurements.\\\bottomrule
\end{tabularx}
\end{table}}

\newcommand{\ServerArchitectureDiagram}{%
\begin{figure}[H]\centering
\begin{tikzpicture}
\node[dBox,fill=bluewash,text width=14.6cm,minimum height=1.15cm] (entry) at (7.55,0)
 {\textbf{HTTPS boundary / Uvicorn + FastAPI}\\Request limits $\rightarrow$ closed Pydantic schemas $\rightarrow$ session lookup $\rightarrow$ resource authorisation};
\node[dBox,text width=4.1cm,minimum height=2.1cm] (care) at (2.3,-2.3)
 {\textbf{Care resources}\\[4pt]Accounts / approved links\\Readings / overview / history\\Consent / events / reminders};
\node[dBox,text width=4.1cm,minimum height=2.1cm] (rules) at (7.55,-2.3)
 {\textbf{Episodes + private inbox}\\[4pt]Source-scoped limit rules\\Persistent alert / recovery\\Per-user delivery + review};
\node[dData,text width=4.1cm,minimum height=2.1cm] (ml) at (12.8,-2.3)
 {\textbf{Synthetic learning}\\[4pt]Classify + save exercises\\Model hash/version check\\Separate SisFall experiment};
\draw[dFlow] (entry.south)--(7.55,-.95)-|(care.north);
\draw[dFlow] (7.55,-.95)--(rules.north);
\draw[dFlow] (7.55,-.95)-|(ml.north);
\node[dBox,fill=light,text width=14.6cm,minimum height=1.25cm] (vault) at (7.55,-4.95)
 {\textbf{Database + Vault / transaction boundary}\\Commit readings and eligible side effects together; unique sample identity protects retries.\\Device diagnostics, synthetic rule input and saved model exercises retain distinct provenance.};
\foreach \n in {care,rules,ml} \draw[dFlow] (\n.south)--(\n.south|-0,-3.85)--(7.55,-3.85)--(vault.north);
\node[dData,text width=14.6cm,minimum height=1.05cm] (db) at (7.55,-6.85)
 {\textbf{SQLite on the self-hosted PC}\\WAL / synchronous FULL / foreign keys; selected sensitive fields protected with AES-256-GCM};
\draw[dFlow] (vault)--(db);
\end{tikzpicture}
\caption{Server runtime responsibilities. Arrows show request/service dependencies and persistence, rather than a single diagnostic pipeline. Physical wearable envelopes are stored as diagnostics; the synthetic model does not classify surface-temperature measurements as core body temperature.}\label{fig:server-architecture}
\end{figure}}

\newcommand{\ServerClassDiagram}{%
\begin{figure}[H]\centering
\begin{tikzpicture}[
 cls/.style={dBox,text width=4.2cm,minimum height=1.8cm},
 dep/.style={dReturn},inherit/.style={-{Triangle[open,length=2.2mm,width=2.4mm]},draw=accent,line width=.85pt}]
\node[cls] (settings) at (2.3,0) {\textbf{Settings}\quad config.py\\[3pt]data\_dir / session\_hours\\load\_key() / initialize\_key()};
\node[cls,fill=bluewash] (database) at (7.55,0) {\textbf{Database}\quad db.py\\[3pt]connect() / transaction()\\audit() / rate\_limit()};
\node[cls] (vault) at (12.8,0) {\textbf{Vault}\quad security.py\\[3pt]seal() / open() / lookup()\\AES-GCM / keyed lookups};
\draw[dep] (database)--(settings);\draw[dep] (database)--(vault);
\node[cls] (schemas) at (2.3,-3.1) {\textbf{ClosedModel}\quad schemas.py\\[3pt]Pydantic BaseModel\\Register / Reading / Event};
\node[cls,fill=bluewash] (api) at (7.55,-3.1) {\textbf{FastAPI route functions}\\[3pt]Session + patient guards\\RequestSecurity middleware};
\node[cls] (health) at (12.8,-3.1) {\textbf{ModelService}\quad ml.py\\[3pt]predict(features, source)\\Saved health-model artefact};
\draw[dep] (api)--(database);\draw[dep] (api)--(schemas);\draw[dep] (api)--(health);
\node[cls] (classification) at (2.3,-6.2) {\textbf{ClassificationRequest}\\routes\_ml.py\\[3pt]Synthetic features + freshness};
\node[cls,fill=bluewash] (services) at (7.55,-6.2) {\textbf{Service modules}\\[3pt]alerts / history / push\\Eligible episodes / rollups / inbox};
\node[cls] (fall) at (12.8,-6.2) {\textbf{FallModelService}\\fall\_ml.py\\[3pt]Independent SisFall model};
\draw[dep] (api)--(services);
\draw[dep] (api.east)--(10.1,-3.1)--(10.1,-4.65)-|(fall.north);
\node[cls] (learning) at (2.3,-9.1) {\textbf{LearningReading}\quad learning.py\\[3pt]sample\_id: UUID};
\draw[inherit] (learning)--(classification);
\node[dNote,text width=9.1cm,anchor=north west] at (5.4,-8.5)
 {\textbf{Legend}\quad Dashed: uses/dependency. Hollow triangle: inheritance.\\[5pt]
  LearningReading extends ClassificationRequest; it does not directly extend ClosedModel. Routes and module functions are shown as functions, not invented controller classes.};
\end{tikzpicture}
\caption{Selected server types and dependencies, derived from the named Python modules. The separate SisFall service remains an experimental model endpoint; it is not a live MPU6050-to-fall alarm pipeline.}\label{fig:server-classes}
\end{figure}}

\newcommand{\DatabaseDiagram}{%
\begin{figure}[H]\centering
\begin{tikzpicture}[
 entity/.style={dBox,text width=3.75cm,minimum height=1.8cm},
 relation/.style={draw=accent,line width=.75pt},
 cardinal/.style={font=\sffamily\fontsize{7.5}{9}\selectfont,text=muted,fill=white,inner sep=1.5pt}]
\node[entity,fill=bluewash] (users) at (2.1,0) {\textbf{users}\\PK id / role\\Encrypted profile\\Argon2id password hash};
\node[entity,fill=bluewash] (patients) at (7.55,0) {\textbf{patients}\\PK/FK id $\rightarrow$ users\\One adult resource};
\node[entity,fill=bluewash] (members) at (13,0) {\textbf{memberships}\\PK elder\_id + caregiver\_id\\Approved adult / caregiver link};
\draw[relation] (users)--node[cardinal,near start,above]{1}node[cardinal,near end,above]{0..1}(patients);
\draw[relation] (patients)--node[cardinal,near start,above]{1}node[cardinal,near end,above]{*}(members);
\draw[relation] (members.north)--(13,1.5)--(2.1,1.5)--(users.north);
\node[cardinal] at (7.55,1.5) {Each membership references 1 caregiver user; a caregiver may hold * memberships};
\node[entity] (sessions) at (2.1,-3.2) {\textbf{sessions + notifications}\\FK user\_id\\Token hash / private inbox\\Expiry / delivery state};
\node[entity] (readings) at (7.55,-3.2) {\textbf{readings + rollups}\\FK elder\_id\\Unique adult/sample identity\\Source + encrypted envelope};
\node[entity] (alerts) at (13,-3.2) {\textbf{alerts + rule state}\\FK elder\_id\\Persistent episode pointer\\Recovery + per-user review};
\draw[relation] (users)--node[cardinal,left]{1 : *}(sessions);
\draw[relation] (patients)--node[cardinal,right]{1 : *}(readings);
\fill[accent] (7.55,-1.5) circle (1pt);
\draw[relation] (7.55,-1.5)--(10.25,-1.5)--(10.25,-3.2)--node[cardinal,above]{1 : *}(alerts.west);
\node[entity] (privacy) at (2.1,-6.5) {\textbf{consents + locations}\\Adult-owned settings\\0..1 consent row per adult\\* expiring location records};
\node[entity] (learning) at (7.55,-6.5) {\textbf{learning\_records}\\FK elder\_id\\Unique adult/sample identity\\Encrypted synthetic exercises};
\node[entity] (care) at (13,-6.5) {\textbf{care + approval records}\\Reminders / occurrences\\Invitations / link requests\\Events / source metadata};
\draw[relation] (7.55,-1.5)--(4.85,-1.5)--(4.85,-5)--(2.1,-5)--(privacy.north);
\draw[relation] (10.25,-3.2)--(10.25,-5.2)--(13,-5.2)--node[cardinal,right]{1 : *}(care.north);
\fill[accent] (10.25,-3.2) circle (1pt);
\draw[relation] (10.25,-5.2)--(7.55,-5.2)--node[cardinal,right]{1 : *}(learning.north);
\fill[accent] (10.25,-5.2) circle (1pt);
\node[dNote,text width=15.1cm,anchor=north west] at (0,-8.05)
 {Grouped ownership map, not the full physical ERD. Patient-owned groups use elder\_id; notifications also reference an adult resource. alert\_acknowledgements uses (alert\_id, user\_id). Operational tables audit and rate\_limits are omitted for clarity.};
\end{tikzpicture}
\caption{Principal SQLite ownership relationships from db.py. Cardinalities describe the database contracts, while grouped records retain their own keys and constraints. Selected fields are encrypted; indexed identifiers and timestamps remain visible metadata.}\label{fig:database-domains}
\end{figure}}

\newcommand{\SecurityProtocolDiagram}{%
\begin{figure}[H]\centering
\begin{tikzpicture}
\node[dHeading,anchor=west] at (0,.65) {Independent trust boundaries};
\node[dBox,fill=bluewash,text width=4.15cm,minimum height=.7cm] at (2.3,0) {\textbf{01 / Account and API}};
\node[dBox,fill=tealwash,text width=4.15cm,minimum height=.7cm] at (7.55,0) {\textbf{02 / Wearable link}};
\node[dBox,fill=bluewash,text width=4.15cm,minimum height=.7cm] at (12.8,0) {\textbf{03 / Private persistence}};
\node[dBox,text width=4.15cm,minimum height=3.8cm,anchor=north] at (2.3,-.8)
 {\textbf{HTTPS + authorisation}\\[6pt]Certificate and host checks\\Argon2id: t=3 / m=64 MiB / p=2\\48 random token bytes\\SHA-256 token hashes\\[6pt]Role + approved relationship\\Expiry and revocation};
\node[dBox,text width=4.15cm,minimum height=3.8cm,anchor=north] at (7.55,-.8)
 {\textbf{SPP + CM1 key possession}\\[6pt]Encrypted Classic Bluetooth\\Private 128-bit board key\\Fresh 128-bit nonces\\HMAC-SHA256 proofs\\[6pt]HELLO / CHALLENGE\\AUTH / READY\\10 s handshake deadline};
\node[dBox,text width=4.15cm,minimum height=3.8cm,anchor=north] at (12.8,-.8)
 {\textbf{Authenticated field storage}\\[6pt]AndroidKeyStore + AES-GCM\\Server Vault: AES-256-GCM\\96-bit random nonce\\Record-bound associated data\\[6pt]Unique sample ID + digest\\Commit before upload ACK};
\node[dData,text width=14.6cm,minimum height=.95cm] at (7.55,-5.65)
 {\textbf{Actual setup order}\quad Adult sign-in $\rightarrow$ private wearable-key setup $\rightarrow$ SPP link and CM1 authentication $\rightarrow$ sensor JSON.\\The columns above group controls; they do not represent a single sequential login protocol.};
\node[dNote,text width=14.8cm,anchor=north west] at (.1,-6.65)
 {Pairing without board input is not proof of resistance to an active pairing relay. Field encryption does not hide SQLite metadata. No passwords, bearer tokens, board keys or proof values appear in the report.};
\end{tikzpicture}
\caption{Implemented security controls across account, radio and storage boundaries. Session authorisation remains independent of Bluetooth key possession. One active phone owns the SPP link; telemetry follows successful CM1 authentication.}\label{fig:security-boundaries}
\end{figure}}


\newcommand{\SequenceDiagram}{%
\begin{figure}[H]\centering
\begin{tikzpicture}[
 font=\sffamily\fontsize{8.5}{11}\selectfont,
 participant/.style={dBox,fill=bluewash,text width=2.35cm,minimum height=.85cm,align=center,inner sep=4pt},
 request/.style={-{Stealth[length=1.7mm]},draw=accent,line width=.8pt},
 reply/.style={-{Stealth[length=1.7mm]},draw=muted,dashed,line width=.65pt},
 msg/.style={dLabel,font=\sffamily\fontsize{8}{10}\selectfont,align=center,fill=white,inner sep=2pt}]
\node[participant] (sender) at (0,0) {\textbf{Sender}\\Synthetic envelope};
\node[participant] (api) at (3.3,0) {\textbf{FastAPI}\\Auth + validation};
\node[participant] (store) at (6.6,0) {\textbf{SQLite}\\Atomic services};
\node[participant] (c1) at (9.9,0) {\textbf{Caregiver A}\\Private inbox};
\node[participant] (c2) at (13.2,0) {\textbf{Caregiver B}\\Private inbox};
\foreach \n in {sender,api,store,c1,c2}
 \draw[stroke,densely dashed,line width=.55pt] (\n.south)--(\n.south|-0,-12.45);
\draw[request] (0,-1.50)--node[msg,above]{POST /readings\\new sample\_id}(3.3,-1.50);
\draw[request] (3.3,-2.30)--node[msg,above]{Check unique key + digest\\Insert encrypted reading}(6.6,-2.30);
\draw[request] (6.6,-2.95)--(7.25,-2.95)--(7.25,-3.60)--(6.6,-3.60);
\node[msg,text width=5.3cm,anchor=west] at (7.7,-3.25)
 {Third eligible abnormal sample:\\one episode; adult, A and B inbox entries};
\draw[reply] (6.6,-4.25)--node[msg,above]{COMMIT; record ID}(3.3,-4.25);
\draw[reply] (3.3,-5.05)--node[msg,above]{201 response lost\\Retry scenario}(0,-5.05);
\draw[request] (0,-5.85)--node[msg,above]{Retry identical envelope\\Same sample\_id}(3.3,-5.85);
\draw[request] (3.3,-6.65)--node[msg,above]{Existing key + identical digest}(6.6,-6.65);
\draw[reply] (6.6,-7.45)--node[msg,above]{Same record; no new alert or fanout}(3.3,-7.45);
\draw[reply] (3.3,-8.25)--node[msg,above]{201; duplicate: true}(0,-8.25);
\draw[request] (9.9,-9.05)--node[msg,above,pos=.48]{A: authenticated long poll}(3.3,-9.05);
\draw[reply] (3.3,-9.85)--node[msg,above,pos=.48]{A's durable inbox item}(9.9,-9.85);
\draw[request] (13.2,-10.65)--node[msg,above,pos=.55]{B: independent authenticated long poll}(3.3,-10.65);
\draw[reply] (3.3,-11.45)--node[msg,above,pos=.55]{B's durable inbox item}(13.2,-11.45);
\draw[request] (9.9,-12.25)--node[msg,above,pos=.48]{A reviews: per-user acknowledgement}(3.3,-12.25);
\node[dNote,text width=14.4cm,anchor=north] at (6.6,-12.85)
 {\tikz{\draw[request](0,0)--(.55,0);}\quad Request / operation\qquad
 \tikz{\draw[reply](0,0)--(.55,0);}\quad Response\\[3pt]
 Changed content with the same sample ID returns 409. Delivery acknowledgement, caregiver review and episode recovery are independent.};
\end{tikzpicture}
\caption{Atomic reception, safe retry and independent recipient fanout. The reading, any eligible rule episode and recipient inbox entries commit together; an identical retry returns the existing record. This contract was verified with loopback clients; simultaneous delivery to two physical phones is not claimed.}\label{fig:sequence}
\end{figure}}

\newcommand{\AlertStateDiagram}{%
\begin{figure}[H]\centering
\begin{tikzpicture}[font=\sffamily\fontsize{8.5}{11}\selectfont,
 state/.style={dBox,text width=3.65cm,minimum height=1.5cm,align=center},
 transition/.style={dFlow},
 edge label/.style={dLabel,font=\sffamily\fontsize{8}{10}\selectfont,align=center,text width=1.25cm}]
\node[state] (idle) at (2.2,0)
 {\textbf{No active episode}\\No active alert pointer\\Per patient and rule};
\node[state,fill=tealwash,draw=teal] (open) at (7.6,0)
 {\textbf{Open episode}\\One persistent alert ID\\Duplicates suppressed};
\node[state] (recovered) at (13,0)
 {\textbf{Recovered episode}\\recovered\_at recorded\\Active pointer cleared};
\draw[transition] (idle.east)--node[edge label,above]{3 present\\abnormal\\estimates}(open.west);
\draw[transition] (open.east)--node[edge label,above]{3 present\\normal\\estimates}(recovered.west);
\draw[transition] (recovered.north)--(13,1.7)--node[dLabel,above]{New three-sample trigger creates a new alert ID}(7.6,1.7)--(open.north);
\draw[transition] (open.south)--(7.6,-1.55)--(5,-1.55)--(5,-.75)--(open.south west);
\node[dLabel,text width=4.8cm,align=center,anchor=north] at (4.65,-1.9)
 {Continued trigger or ineligible input:\\keep the current episode open};
\node[dBox,fill=bluewash,text width=5.4cm,minimum height=1.3cm] (review) at (12.1,-2.5)
 {\textbf{Independent caregiver review}\\Record (alert\_id, user\_id)\\Review does not close or recover an episode};
\draw[dReturn] (open.south east)--(9.65,-1.25)-|(review.north);
\node[dNote,text width=14.5cm,anchor=north] at (7.6,-3.65)
 {\textbf{Eligibility gate:} the latest three readings must share patient and source, be fresh within two minutes, correctly ordered and 1--30 seconds apart. Optical rules use present numerical estimates with contact; quality metadata is retained. Missing contact, stale data or absent estimates cannot prove recovery. Device and synthetic episodes are independent.};
\node[dNote,anchor=north west] at (.05,-5.15)
 {\tikz{\draw[dFlow](0,0)--(.6,0);}\quad Episode transition\qquad
 \tikz{\draw[dReturn](0,0)--(.6,0);}\quad Independent review action};
\end{tikzpicture}
\caption{Implemented educational episode lifecycle in alerts.py. A recovered alert remains in history; only its active pointer is cleared. ``Recovered'' refers to a configured rule episode, not a measured health outcome.}\label{fig:alert-state}
\end{figure}}

\newcommand{\MobileLayersDiagram}{%
\begin{figure}[H]\centering
\begin{tikzpicture}[font=\sffamily\fontsize{8.5}{11}\selectfont,
 module/.style={dBox,text width=4.1cm,minimum height=1.8cm}]
\node[dBox,fill=bluewash,text width=14.4cm,minimum height=1.2cm] (activity) at (7.6,0)
 {\textbf{PRESENTATION AND CONTROLLER}\\MainActivity + LearningActivity: native Views and screen state\\MetricHistoryChartView draws aggregates; executor results return to the main thread.};
\node[module] (helpers) at (2.3,-3)
 {\textbf{Logic and state helpers}\\ReadingSummary\\LiveSampleFreshness\\RelayPolicy; AlertInbox};
\node[module] (api) at (7.6,-3)
 {\textbf{API and account boundary}\\ApiClient: HTTPS + JSON\\ApiSession: identity + session\\Authorisation remains on server};
\node[module] (store) at (12.9,-3)
 {\textbf{Private persistence}\\SecureStore: AES-GCM files\\AndroidKeyStore key\\Account-bound state};
\draw[dFlow] (activity.south west)--(helpers.north);
\draw[dFlow] (activity.south)--(api.north);
\draw[dFlow] (activity.south east)--(store.north);
\node[module,fill=tealwash] (transport) at (2.3,-6.55)
 {\textbf{Device reception}\\LiveMonitorService\\SensorTransport interface\\SppSensorClient\\BleSensorClient};
\node[module,fill=tealwash] (relay) at (7.6,-6.55)
 {\textbf{Durable upload worker}\\ServerRelay\\Persist before sending; stable UUID\\Bounded queue; retry policy};
\node[module,fill=tealwash] (notices) at (12.9,-6.55)
 {\textbf{Private inbox service}\\NotificationChannelService\\NotificationSync\\NotificationCursor};
\draw[dFlow] (transport.east)--node[dLabel,above]{snapshot}(relay.west);
\draw[dFlow] (relay.north)--node[dLabel,right]{HTTPS}(api.south);
\draw[dFlow] (notices.north)--node[dLabel,right]{cursor}(store.south);
% Notification delivery also depends on the API, rather than on storage alone.
\draw[dFlow] (notices.west)--(10.25,-6.55)--(10.25,-3)--(api.east);
\node[dLabel,rotate=90,fill=white] at (10.25,-4.75) {HTTPS long poll};
\node[dNote,text width=14.8cm,anchor=north] at (7.6,-8.05)
 {\tikz{\draw[dFlow](0,0)--(.65,0);}\quad Selected runtime calls and data transfers. Each row groups responsibilities; it is not a strict dependency prohibition.\\[3pt]
 Activities combine View and controller work in a layered, MVC-inspired design. BLE is a separate client option; the current ESP32 image uses SPP. Metric history and saved synthetic inference both have native screens.};
\end{tikzpicture}
\caption{Actual Android organisation. Reusable transport, upload, notification, HTTPS and state helpers support the activities. Source separation and JVM tests do not alone establish Android lifecycle, accessibility, Bluetooth or background-delivery acceptance.}\label{fig:mobile-layers}
\end{figure}}
